\begin{tikzpicture}[font=\sffamily]
  % Mapa de Empatía clásico alrededor de la figura central del cliente
  \node[font=\sffamily\bfseries\small, ceindigo] at (0, 4.45)
    {MAPA DE EMPATÍA DEL SEGMENTO OBJETIVO (OSTERWALDER / XPLANE)};

  % Caja rectangular exterior
  \draw[ceindigo, line width=1.1pt, rounded corners=4pt] (-6.8, -3.95) rectangle (6.8, 3.95);
  % Divisiones diagonales hacia el círculo central (centro (0,0.85), radio 1.05)
  \draw[ceindigo!70, line width=0.8pt] (-6.8, 3.95) -- (-0.955, 1.286);
  \draw[ceindigo!70, line width=0.8pt] (6.8, 3.95)  -- (0.955, 1.286);
  \draw[ceindigo!70, line width=0.8pt] (-6.8, -1.65) -- (-0.986, 0.488);
  \draw[ceindigo!70, line width=0.8pt] (6.8, -1.65)  -- (0.986, 0.488);
  % Separador inferior Esfuerzos vs Resultados
  \draw[ceindigo!80, line width=0.9pt] (-6.8, -1.65) -- (6.8, -1.65);
  \draw[ceindigo!80, line width=0.9pt] (0, -3.95) -- (0, -1.65);

  % Núcleo central del Cliente
  \fill[cepanel] (0, 0.85) circle (1.05cm);
  \draw[ceindigo, line width=1.2pt] (0, 0.85) circle (1.05cm);
  \node[align=center, font=\sffamily\bfseries\scriptsize, ceindigo] at (0, 0.85)
    {PERSONA /\\CLIENTE\\OBJETIVO};

  % Cuadrante superior: Qué piensa y siente
  \node[text width=5.6cm, align=center, font=\sffamily\scriptsize] at (0, 3.05) {
    \textcolor{cecrimson}{\bfseries 1. ¿QUÉ PIENSA Y SIENTE?}\\[2pt]
    \color{ceslate} Preocupaciones reales, aspiraciones íntimas, prioridades y emociones que no siempre expresa en voz alta.
  };

  % Izquierda: Qué ve
  \node[text width=3.8cm, align=left, font=\sffamily\scriptsize] at (-4.35, 0.65) {
    \textcolor{ceteal!90!black}{\bfseries 2. ¿QUÉ VE?}\\[2pt]
    \color{ceslate} Entorno próximo, oferta de competidores, qué hacen colegas y amigos, tendencias en su sector.
  };

  % Derecha: Qué oye
  \node[text width=3.8cm, align=left, font=\sffamily\scriptsize] at (4.35, 0.65) {
    \textcolor{ceteal!90!black}{\bfseries 3. ¿QUÉ OYE?}\\[2pt]
    \color{ceslate} Qué dicen jefes, clientes, familiares y prescriptores; medios o canales que influyen en su criterio.
  };

  % Centro-abajo (entre círculo y línea inferior): Qué dice y hace
  \node[text width=7.2cm, align=center, font=\sffamily\scriptsize] at (0, -0.96) {
    \textcolor{ceindigo}{\bfseries 4. ¿QUÉ DICE Y HACE?}\\[1pt]
    \color{ceslate} Conducta observable, hábitos diarios, contradicciones entre discurso y acción real.
  };

  % Inferior izquierda: Esfuerzos / Miedos
  \fill[cepanelcrimson, opacity=0.7] (-6.78, -3.93) rectangle (-0.02, -1.67);
  \node[text width=6.1cm, align=left, font=\sffamily\scriptsize] at (-3.4, -2.8) {
    \textcolor{cecrimson}{\bfseries 5. ESFUERZOS (Miedos, frustraciones, obstáculos)}\\[2pt]
    \color{ceslate} Qué riesgos teme al cambiar de proveedor, qué barreras de tiempo/coste le frenan hoy.
  };

  % Inferior derecha: Resultados / Beneficios
  \fill[cepanelteal, opacity=0.7] (0.02, -3.93) rectangle (6.78, -1.67);
  \node[text width=6.1cm, align=left, font=\sffamily\scriptsize] at (3.4, -2.8) {
    \textcolor{ceteal!90!black}{\bfseries 6. RESULTADOS (Deseos, medidas de éxito)}\\[2pt]
    \color{ceslate} Cómo mide el éxito real en su trabajo o vida; qué logro concretaría su compra.
  };
\end{tikzpicture}